% 本文件由 LaTeX 排版 Agent 根据有效 translation.json 自动生成。
% author=Athenagoras; work_id=0205; title=为基督徒申辩

\chapter{为基督徒申辩}

% structure: p0001 source=h1 target=omitted reason=page-title-already-rendered-by-parent

\Signature{\Emph{雅典人阿特纳戈拉：哲学家与基督徒}}\par

\Emph{致皇帝\Person{马可·奥勒留·安东尼努斯}与\Person{卢修斯·奥勒留·康茂德}，\Place{亚美尼亚}与\Place{萨尔马提亚}的征服者，且最重要的，哲学家。}\par

% structure: p0004 source=h2 target=section reason=relative-heading-rank; base=h2

\section{第一章 基督徒遭受的不公}

在你们帝国里，最伟大的君主们，不同的民族有不同的习俗和法律；没有人因法律或惩罚的恐惧而受阻，无法遵从他祖先的惯例，无论这些惯例多么可笑。\Place{伊利昂}的一位公民称\Person{赫克托耳}为神，并将\Person{海伦}当作\Person{阿德拉斯忒娅}来予以神性尊崇。\Person{拉栖代梦人}尊奉\Person{阿伽门农}为\Person{宙斯}，尊奉\Person{廷达瑞俄斯}的女儿\Person{菲洛诺厄}；\Place{特涅多斯}人敬拜\Person{坦尼斯}。\Person{雅典人}向\Person{厄瑞克透斯}献祭，视其为\Person{波塞冬}。\Person{雅典人}还举行宗教仪式并庆祝密仪，尊崇\Person{阿格劳洛斯}与\Person{潘德洛索斯}，这两位女子曾因打开箱子而被判犯有不敬之罪。简言之，在每一个民族和人民中，人们随意献祭，庆祝他们所喜欢的任何密仪。\Person{埃及人}甚至将猫、鳄鱼、蛇、毒蛇和狗列为神祇。对于所有这些，你们和法律都允许如此行事，一方面认为不信任何神是不敬且邪恶的，另一方面认为每个人必须敬拜他所偏好的神，以便因对神明的畏惧而使人们远离恶行。但是，为什么——请不要像大众那样，被传闻引入歧途——为什么仅仅一个名字就令你们憎恶呢？名字本身不应招致仇恨：应受惩罚与刑罚的是不义的行为。因此，我们钦佩你们的温和与仁慈，以及你们对每个人的和平与仁爱倾向，个人享有平等权利；各城市按等级分享同等荣誉；在你们睿智的治理下，整个帝国享有深度的和平。但对于我们这些被称为基督徒的人，你们却未以同样方式关怀；尽管我们没有犯下任何罪行——不，正如下文所述，我们所有人中，对\OriginalTerm{天主}{Deity}和对你们的政府最为虔诚与正直——你们却允许我们遭受骚扰、掠夺与迫害，众人仅因我们的名字就向我们开战。因此，我们斗胆向你们陈述我们的案情——你们将从这篇申辩中得知，我们遭受的苦难是不公的，且违背一切法律与理性——我们恳求你们也能对我们有所考虑，使我们终于不再因虚假控告者的唆使而被屠杀。因为迫害者对我们课处的罚金并非仅仅针对我们的财产，他们的侮辱并非仅仅针对我们的名誉，他们加于我们的损害也并非针对我们其他更重大的利益。这些事我们视如敝履，尽管对大众来说它们似乎至关重要；因为我们已学到，不仅不以拳还拳，不与掠夺抢劫我们的人打官司，而且对那些打我们这边脸的人，要把另一边脸也转给他；对那些拿走我们内衣的人，连外衣也要给他。\BibleRef{玛 5:39-40}\par

% structure: p0006 source=h2 target=section reason=relative-heading-rank; base=h2

\section{第二章 要求在被控告时受同等待遇}

诚然，若有人能指证我们犯有何罪，无论大小，我们并不请求免于刑罚，反已准备好承受最严厉、最无情的惩处。但若控告仅关乎我们的名号——无可否认，迄今关于我们的传说，无非建立在普遍而漫无鉴别的流言之上，且从未有基督徒被证实犯有罪行——那么，卓越、仁慈且博学的君主们，你们便当以法律除去这种侮辱性的待遇，以便正如举世之民及邦国皆蒙受你们的恩泽，我们也能对你们心怀感激，欢欣于不再成为诬告的牺牲品。因为你们的公义不容许这样的事：他人被控犯罪时，未证实罪状即不加以刑罚，而在我们的情形中，我们所背负的名号竟比审讯中提出的证据更具力量，法官们不去查问被告是否犯有何罪，反倒在名号上发泄侮辱，仿佛名号本身即是一桩罪行。然而，任何名号本身皆不被视为善或恶；名号之显为恶或善，端视其背后所对应的行为是恶是善。对此，你们自己亦深知，因为你们深谙哲学及一切学问。正因如此，那些被带到你们面前受审的人，即便被控以最严重的罪名，也毫不畏惧，因他们知道你们必会查问其过往的生活，而不会被空洞的名号所左右，亦不会因诉状中所陈的指控为虚假而受影响：他们就判罪或开释之公正而欣然接受判决。是以，我们为自己索求那普遍承认的共同权利：我们不应因被称为基督徒而遭憎恨与惩罚（名号与我们之为恶人何干？），而是当按所控之任何罪名受审，若驳斥了它们，即被开释，若证实有罪，才受惩罚——非为那名号（因基督徒若非假充我们的教义，没有一个是恶人），而是为所犯的恶行。我们看到哲学家们正是如此受审的。他们在受审之前，法官绝不因其学识或技艺而视之为善或恶；若证实犯有恶行，便受惩罚，而哲学并不因此蒙受污名（因他是不按正途修习哲学，故为恶人，而学识本身无可指摘），若他驳倒了虚假的控告，则被开释。那么，也当以这同等的公义对待我们。当查究被控者的生活，而名号却应免受一切归咎。在我申辩之初，须恳请诸位杰出的皇帝，容我以不偏不倚之心垂听：请勿随从普遍的、失却理性的议论而预断此案，却将你们的求知之欲与热爱真理之心也应用于审查我们的教义。这样，你们一方面不致因无知而错判，我们方面则因驳斥了出自不加鉴别的群众传闻的种种指控，而不再遭受攻击。\par

% structure: p0008 source=h2 target=section reason=relative-heading-rank; base=h2

\section{第三章 对基督徒的控告。}

受人指控的共有三件事：无神论、\OriginalTerm{提厄斯忒斯的宴席}{Thyestean feasts}、\OriginalTerm{俄狄浦斯式的交合}{Œdipodean intercourse}。然而，若这些控告属实，便不必饶恕任何人：径直着手处置我们的罪行吧；若发现有任何基督徒过着如野兽般的生活，就连同我们的妻子儿女，根除灭绝我们。然而，即便是野兽也不尝同类之血肉；它们依循自然法则交配，且只在合宜的季节，并非纯粹出于淫乱；它们也认得那施恩于己者。因此，若有人比野兽更加凶残，那么，什么刑罚能算作与他所犯的暴行相称呢？但是，如果这些事不过是无聊的故事和虚妄的毁谤，其根源在于德行与恶行天然对立，而相反之物循着神圣的律法彼此交战（你们自己便是见证，因你们禁止人向我们告密），那么，就有待你们探究我们的生活、我们的见解、我们对你们及你们的家族与政府的忠心服从，并最终赐予我们与迫害我们之人同等的权利（我们别无所求）。因为那时我们必能胜过他们，正如现今一样，为了真理，我们毫不犹豫地舍弃自己的性命。\par

% structure: p0010 source=h2 target=section reason=relative-heading-rank; base=h2

\section{第四章 基督徒并非无神论者，而是承认唯一的天主。}

首先，关于我们为无神论者的指责——我将逐一回应各项控告，以免因无辞以对指控者而遭嗤笑——雅典人将\Person{迪亚戈拉斯}判为无神论者，是有道理的，因为他不仅泄露了\OriginalTerm{俄尔甫斯教义}{Orphic doctrine}，公开了\Place{厄琉西斯}及\OriginalTerm{卡比里}{Cabiri}的秘仪，又将\Person{赫拉克勒斯}的木像劈来煮他的萝卜，更公然宣称根本不存在任何神。然而，我们却是将天主与物质分别开来，并教导说物质是一回事，天主是另一回事，两者间隔着巨大的鸿沟（因为天主是非受造的、永恒的，只能由理智与理性所洞见，而物质则是受造的、可朽坏的），若将无神论之名加诸我们，岂不荒谬？倘若我们的想法与迪亚戈拉斯相同，而面对的却是如此激发虔敬的种种动力——那既定的秩序、宇宙的和谐、广大、色泽、形态、世界的安排——那么，我们被斥为不虔敬的名声，以及我们遭受如此迫害的原因，便理当归咎于我们自己。然而，既然我们的教义承认唯一的天主，即这宇宙的造物主，祂自身是非受造的（因为\Quote{是}者不生，\Quote{不是}者方生），却借由出自祂的\OriginalTerm{圣言}{Logos}创造了万物，因而我们在两方面都蒙受了不合理的对待：既遭诽谤，又受迫害。\par

% structure: p0012 source=h2 target=section reason=relative-heading-rank; base=h2

\section{第五章 诗人对天主的唯一性的见证。}

诗人们与哲学家们探寻关于天主的事，并未被定以无神论者的罪名。\Person{欧里庇得斯}论及那些据庸常之见蒙昧地被称为神者，迟疑地说道：——\par

\Quote{若\Person{宙斯}果真在天上为王，\newline{}就不该将灾祸降于义人。}\par

但论及那位凭理智可确知者，他却以理智明断地给出了意见，如下：——\par

\Quote{你可看见在高处，那以湿润双臂\newline{}拥抱无垠苍穹与大地的？\newline{}你要以他为\Person{宙斯}，视他为天主。}\par

因为，对于这些所谓的众神，他既未看到任何真实存在作为其基础而被冠以通常的名字（例如 \Quote{宙斯}：\Quote{宙斯是谁我不知，只是听人说}），也未看到有任何名字被赋予实际存在的事物（因为对那些没有真实存在作为基础者，名字又有何用？）；但他却借着祂的工程看见了祂，以洞察无形之物的眼目，默观了在空气、以太和大地中显现之物。因此，那创造万物、并以祂的圣神治理万物的祂，他推断就是天主；\Person{索福克勒斯}与他一致，说：——\par

\Quote{只有一位天主，确实只有一位，\newline{}创造了高天和辽阔大地的那一位。}\par

[\Person{欧里庇得斯}] 论及天主的本性，这本性使祂的工程充满美，并教导天主应在何处，以及祂须是唯一。\par

% structure: p0020 source=h2 target=section reason=relative-heading-rank; base=h2

\section{第六章 哲学家关于独一天主的意见。}

\Person{菲洛劳斯}也曾说，万物包含在天主内如同在堡垒中，他教导天主是唯一的，且高于物质。\Person{吕西斯}和\Person{奥普西莫斯}这样定义天主：一位说天主是不可言喻的数字，另一位说天主是最大数超过次大数的优越。因此，既然按毕达哥拉斯学派，十是最大数，即\OriginalTerm{四元体}{Tetractys}，包含一切算术与和谐原理，而九紧随其后，那么天主就是一个单元——即，唯一。因为最大数超过次大数一。还有柏拉图和亚里士多德——并非我要逐一叙述哲学家关于天主的全部言论，好像我要展示他们意见的完整摘要似的；因为我知道，正如你在理智和统治权柄方面超越所有人，同样，你在精确通晓一切学问方面也超越他们，你以比那些专攻一门的人更大的成就研习各个分支。然而，由于若非引述名字，便不可能证明并非只有我们将天主概念限定为单一，所以我冒昧列举这些意见。于是，柏拉图说：\Quote{要发现这宇宙的创造者和父是困难的；即使发现了，也不可能向所有人说明，} 他认为有一位非受造的永恒天主。如果他承认还有其他者，如日、月、星辰，他也承认它们是受造的：\Quote{神，众神的后裔，我是它们的创造者，及那些若非我意愿便不可分解的工程之父；但凡组合之物皆可分解。} 因此，如果柏拉图因认为有一位非受造的宇宙创造者天主而不算是无神论者，那么我们也不算是无神论者，我们承认并坚定持守，那位借圣言创造万物、并借祂的圣神维系万有的，就是天主。亚里士多德及其追随者承认一位可视为某种复合\OriginalTerm{活物}{\GreekText{ζῶον}}的存在，说天主由灵魂和身体构成，认为祂的身体是苍穹空间、行星及恒星天球，作圆周运动；而祂的灵魂是主宰身体运动的理性，自身不动，却为其他运动之因。斯多葛学派亦如此，尽管他们为了适应物质的变化而使用各种名称——他们说物质被天主的灵渗透——他们在名义上使神性增多，但实际上认为天主是唯一的。因为，如果天主是有计划地推动世界万物产生的艺术之火，自身包含一切种子原理，万物依命运藉此产生，而如果祂的灵遍布整个世界，那么按照他们，天主便是唯一的，因物质的炽热部分（\OriginalTerm{沸腾者}{\GreekText{τὄ} \GreekText{ζέον}}）而称为\Person{宙斯}，因\OriginalTerm{空气}{\GreekText{ὁ} \GreekText{ἀήρ}}而称为\Person{赫拉}，并因祂所渗透的物质的特定部分而被称为其他名称。\par

% structure: p0022 source=h2 target=section reason=relative-heading-rank; base=h2

\section{第七章 基督徒关于天主学说的优越性。}

既然，几乎所有人都承认天主性的唯一，即便出于勉强，当他们论及宇宙的第一原理时，而我们也同样断言安排这宇宙者是天主——为何他们可以任意谈论和书写关于天主性的事而不受惩罚，却有法律加于我们，尽管我们能以符合真理的证据和理性证明我们所领会并合理相信的，即只有一位天主？因为诗人和哲学家们，一如对其他课题，对此课题亦采取臆测之途，因与来自天主的灵感有某种亲和，各人藉自己的灵魂被推动，试图寻获并领会真理；但他们未显出足以完全领会它，因为他们认为宜于不从天主学习关于天主的事，而是各自从自己学习；于是他们各人对天主、物质、形式和世界得出自己的结论。但我们有先知为我们所领会和相信之事作见证，他们是宣讲天主及天主之事的人，由圣神引导。你也会承认，既然你在理智和对\OriginalTerm{真天主}{\GreekText{τὸ} \GreekText{ὄντως} \GreekText{θεῖον}}的虔敬上超越所有人，若我们停止相信来自天主的圣神——祂像乐器一样推动先知的口——而听从于纯粹属人的意见，那是非理性的。\par

% structure: p0024 source=h2 target=section reason=relative-heading-rank; base=h2

\section{第八章 多神论的荒谬。}

论到那从起初便有一位天主，即这宇宙的\OriginalTerm{造物主}{Maker}，的教义，请你如此思考，好使你也能了解我们信仰的论证依据。如果从起初有两位或更多位神，他们要么在同一地方，要么各自分处不同地方。在同一地方不可能。因为，如果是神，就不会相同；正因他们是\OriginalTerm{未受造的}{uncreated}，他们便不相似：因为受造之物与其原型相似；但非受造者不相似，既非由某一位所生，也非按某一位的模子而造。手、眼和脚是同一个身体的肢体，合起来构成一个人：天主是否如此是一位呢？事实上，苏格拉底是组合而成且可分的，正因他是受造的且可朽坏；而天主则是未受造的，\OriginalTerm{无欲情的}{impassible}，\OriginalTerm{不可分的}{indivisible}——因此不由部分组合。然而，如果相反，他们各自分开存在，那么，既然创造世界的那一位超越受造物，且在其所造并安排的事物之上，那第二位或其余各位能在哪里呢？因为，如果世界被造成球形，并限定在天穹的圈内，而世界的创造者超越受造物，以祂的\OriginalTerm{天主眷顾}{providential care}治理它们，那么，那里有地方给第二位神或其余诸神呢？因为他不在世界之内，因为世界属于另一位；也不在世界周围，因为创造世界的天主在其上。但如果他既不在世界内也不在世界周围（因为周遭全被这一位所占据），那么他在哪里呢？他是在世界和[第一位]天主之上吗？在另一个世界或另一个世界周围吗？但如果他在另一个世界或其周围，那他便不在我们之上，因为他并不治理世界；他的能力也不伟大，因为他存在于有限的空间内。但如果他既不在另一个世界（因为万物已被另一位充满），也不在另一个周围（因为万物已被另一位占据），他显然根本不存在，因为没有地方可容他存在。或者，他有另一位拥有世界，而他则在世界造物主之上，却又既不在世界内也不在世界周围，他能做什么呢？那么，有某个其它地方吗？然而，天主及属于天主的一切却在他之上。还有，既然另一位充满了世界以上的区域，何来地方？或许他施加天主眷顾？[绝非如此。] 然而，除非他这样做，他便一事无成。因此，如果他既不行事也不施天主眷顾，且没有其它地方可让他存在，那么我们所说的这位存在者就是从起初唯一的真天主，也是世界的独一造物主。\par

% structure: p0026 source=h2 target=section reason=relative-heading-rank; base=h2

\section{第九章 先知们的见证。}

如果我们只满足于提出这些理由，我们的教义可能会被某些人视为人间的学说。然而，由于先知们的声音证实了我们的论证——我想，以你如此渴求知识、学养卓越，也绝不会不知道\Person{梅瑟}的著作，以及\Person{依撒意亚}、\Person{耶肋米亚}和其他先知的书卷；他们在\OriginalTerm{圣神}{Divine Spirit}的推动下，\OriginalTerm{神魂超拔}{ecstasy}，超越了理智的自然运作，说出所受的灵感之言，圣神利用他们，如同吹笛者向笛内吹气——那么，这些人说了什么呢？\Quote{上主是我们的天主；没有谁能与祂相比。}又说：\BibleQuote{我是天主，元始和终结，除我以外，没有别的神。}\BibleRef{依 44:6}同样：\BibleQuote{在我以前，没有别的神；在我以后，也绝不会有；我是天主，除我以外没有别的神。}\BibleRef{依 43:10}至于祂的伟大：\BibleQuote{天是我的宝座，地是我脚的脚凳：你们要为我建造怎样的殿宇？哪里是我安息之所？}\BibleRef{依 66:1}但我留待你亲自查阅这些书卷时，仔细省察其中所含的预言，好使你能以合适的理由为我们的信仰辩护，对抗所加于我们的辱骂。\par

% structure: p0028 source=h2 target=section reason=relative-heading-rank; base=h2

\section{第十章 基督徒敬拜圣父、圣子及圣神。}

所以我因我们承认一位未受造、永恒、无形、不受苦难、不可理解、无限的天主，这位天主只能由理智和理性把握，为光、美、灵和不可言喻的德能所环绕，宇宙是藉着祂的圣言而受造、井然有序并得以存立——我已充分证明。[我说\Quote{祂的圣言}]，因为我们还承认一位圣子。没有人应觉得天主有圣子是可笑的。因为虽然诗人们在他们虚构的故事中把神描绘得不比人强，我们关于圣父或圣子的思想方式，却与他们不同。但圣子是天主圣父的圣言，在理念上和在实施上都是；因为万物是按祂的模范并藉着祂而造成的，圣父与圣子原为一体。且圣子在圣父内，圣父在圣子内，在合一与灵性能力上，圣父的\OriginalTerm{理智}{\GreekText{νοῦς}}与\OriginalTerm{理性}{\GreekText{λόγος}}就是圣子。然而，如果出于你们超绝的智慧，你们想查问圣子是什么意思，我将简略地说明：祂是圣父的初产物，并非如同被造而生（因为从起初，天主——那永恒的\OriginalTerm{理智}{\GreekText{νοῦς}}，就在自身内拥有圣言，自永恒起就具有\OriginalTerm{逻各斯的}{\GreekText{λογικός}}性质）；而是作为一切物质事物的理念与驱动力量而发出，那些物质事物原本如同一团没有属性的自然、一片没有活力的土地，较重的颗粒与较轻的混在一起。预言之灵也同意我们的说法。祂说：\BibleQuote{上主造了我，作他造化的起始，在他一切作为之初。}\BibleRef{箴 8:22} 我们断言，那运行在先知中的圣神本身也是从天主流出的，从祂发出，又像阳光一样返回。那么，听到有人谈论圣父、圣子与圣神，并宣告他们既在合一中拥有权能又在次序上有区别，却被称作无神论者，谁不感到惊讶呢？我们在神圣本性方面的教导也并非仅限于这些要点；我们还承认有众多天使和仆役，是造物主和世界的设计者天主藉着祂的圣言分配并安置到各自的岗位上，去掌管元素、诸天、世界及其中的万物，并维系它们全体的美好秩序。\par

% structure: p0030 source=h2 target=section reason=relative-heading-rank; base=h2

\section{基督徒的伦理教导驳斥了施加于他们的指控。}

如果我详述我们的教义细节，请勿惊讶。这是为了使你们不被世俗和非理性的意见所左右，而能将真理清楚地置于你们眼前。因为将我们所持守的见解本身——这些见解并非出于人，而是由天主说出并教导的——展现出来，我们就能说服你们不把我们看作无神论者。那么，我们自幼所受的教导是什么呢？\BibleQuote{我告诉你们：要爱你们的敌人，祝福诅咒你们的人，为迫害你们的人祈祷，这样你们就可以做你们在天之父的儿女；因为他使太阳上升，光照恶人，也光照善人；降雨给义人，也给不义的人。}\BibleRef{路 6:27} 请允许我在此放胆扬声，以洪亮可闻的呼声，在爱好哲学的君王面前辩护。因为那些解析三段论、厘清歧义、解释词源的人，那些教导同音异义词和同义词、范畴和公理、何为主词何为谓词的人，那些向门徒许诺通过这些教导能获得幸福的人：他们中有谁洁净了自己的灵魂，以至于不是恨敌人，而是爱他们；不是辱骂中伤他们的人（仅仅是避免咒骂，本身就彰显了不凡的忍耐），而是祝福他们；并为那些图谋害他们性命的人祈祷？相反，他们怀着恶意不断地巧妙探究他们技艺的奥秘，并总是致力于做坏事，以言辞之术而非行为的表现为职业和专业。但在我们当中，你们会发现没有受过教育的人、工匠和老妇，他们若不能用言语证明我们的教义之效，却以他们的行为显明了他们因坚信其真理而获得的益处：他们不背诵演说词，而是展示善行；被打而不还手；被偷盗而不诉诸法律；他们向求他们的人施舍，并爱邻如己。\par

% structure: p0032 source=h2 target=section reason=relative-heading-rank; base=h2

\section{无神论指控的必然荒谬}

那么，除非我们相信有一位天主统御人类，我们难道会如此涤除邪恶吗？绝不会。然而，因为我们深信，对于今生的一切，我们必将向创造我们和世界的天主交账，所以我们采取一种节制、仁爱、普遍受人蔑视的生活方式，相信即便我们的性命被夺去，我们在此所受的大恶，与我们从大审判者那里因温良、仁爱、节制的生活而将获得的相比，也是微不足道的。\Person{柏拉图}曾说，\Person{米诺斯}和\Person{拉达曼提斯}将审判并惩罚恶人；但我们说，即使一个人是\Person{米诺斯}或\Person{拉达曼提斯}本人，或是他们的父亲，他也无法逃脱天主的审判。那么，那些认为人生仅在于此——\BibleQuote{我们吃喝吧！因为明天就要死了}——并视死亡为沉睡与遗忘（\Quote{睡眠和死亡，孪生兄弟}）的人，竟被视为虔诚吗？而我们这些估量今生价值甚微，唯独因认识天主和祂的\OriginalTerm{圣言}{Logos}，认识圣子与圣父如何为一，圣父与圣子如何共融，圣神是什么，这三位——圣神、圣子、圣父——如何合一并在合一中区分；且知道我们所盼望的生命远比言语所能描述的更为美好，只要我们保持纯净，毫无过犯，从而抵达那里；此外，我们的仁爱推及如此地步，以至于我们不仅爱我们的朋友（正如祂说：\BibleQuote{因为如果你们爱那些爱你们的人，借钱给那些借钱给你们的人，你们还有什么赏报呢？}）——我说，当我们有这样的品格，过这样的生活，以期最终免于定罪时，我们难道不算是虔诚的吗？然而，这些不过是大事中的小事，许多事中的几件，免得我们再耗费您的耐心；因为品尝蜂蜜与乳清的人，只凭少量便能判断整体是否美好。\par

% structure: p0034 source=h2 target=section reason=relative-heading-rank; base=h2

\section{第十三章 基督徒为何不献祭}

但是，大多数控告我们为无神论者的人——之所以如此，是因为他们对天主是什么连最模糊的概念也没有，他们愚钝，对属自然和属神的事物全然无知，并且以祭献的规条来衡量虔诚——控告我们不承认各个城邦所崇拜的那些神明。诸位皇帝，请留意以下两方面的问题。首先，关于我们不献祭：这宇宙的塑造者和父不需要血液，不需要全燔祭的馨香，也不需要花朵和香料的芬芳，因为祂本身就是完美的芬芳，既不欠缺内在的什么，也不欠缺外在的什么；然而，对祂而言，最崇高的祭献，就是我们认识那位铺张穹苍、安立大地如中心、聚集海水成海洋、分开光明与黑暗、以星辰装饰天空、使大地生出各从其类的种子、创造动物并塑造人类的天主。当我们怀着这样的信念，即天主是万有的塑造者，以知识和治理的智慧保存万有并管理一切，我们向祂\BibleQuote{举起圣洁的手}，祂又何需百牛大祭呢？\par

\Quote{因为当凡人犯罪或未能\newline{}行义时，借祭献和祈祷，\newline{}奠酒与全燔祭，可获得安抚。}\par

至于全燔祭，天主并不需要，我又与它何干呢？——然而，我们理当献上不流血的祭献和\BibleQuote{我们理智的敬礼}。\par

% structure: p0038 source=h2 target=section reason=relative-heading-rank; base=h2

\section{第十四章 控告基督徒者的自相矛盾}

然后，关于另一个指控，即我们不向各城邦所拜的相同神明祈祷和信仰，这是极为愚蠢的。为什么，正是那些指控我们为无神论、因为我们不承认与他们一样的神明的人，他们自己在神明问题上也互不一致。雅典人以\Person{刻琉斯}和\Person{墨塔尼拉}为神；拉栖第梦人以\Person{墨涅拉俄斯}为神，向他献祭、举行庆典，而\Place{伊利昂}人却连他的名字都不堪入耳，却崇拜\Person{赫克托耳}。凯奥斯人崇拜\Person{阿里斯泰俄斯}，视之为\Person{宙斯}和\Person{阿波罗}；塔索斯人崇拜\Person{忒阿革涅斯}，这人曾在奥林匹克竞技中杀人；萨摩斯人崇拜\Person{吕山德}，尽管他犯下所有屠杀和罪行；\Person{阿尔克曼}和\Person{赫西俄德}崇拜\Person{美狄亚}，奇里乞亚人崇拜\Person{尼俄柏}；西西里人崇拜\Person{布塔基得斯之子腓力}；阿马苏斯人崇拜\Person{奥涅西洛斯}；迦太基人崇拜\Person{哈米尔卡}。时间不允许我一一列举。那么，他们自己在神明问题上互不一致，又为何控告我们与他们不一致呢？再看埃及人中流行的做法：岂非极为可笑？他们在重大节日里在神庙中如同为死者捶胸哀悼，又将之作为神明献祭；这并不奇怪，因为他们视牲畜为神，牲畜死后他们剃发，将其埋葬在神庙中，并公开哀悼。那么，如果我们因没有实行符合他们方式的敬神而被视为不虔，那么所有城邦和所有民族都犯了不虔之罪，因为他们并不都承认相同的神明。\par

% structure: p0040 source=h2 target=section reason=relative-heading-rank; base=h2

\section{第十五章 基督徒将天主与物质区分}

但即使承认他们也承认同样的道理，那又如何？因为大众不能区分物质与天主，也不能看见二者之间有多么巨大的鸿沟，他们向用物质造成的偶像祈祷；难道我们，因我们能分辨并区分非受造的与受造的、存在者与不存在者、理智所领悟的与感官所感知的，且给予每一类相称的名称——我们就该去叩拜偶像吗？若物质与天主真是一回事，同一个事物的两个名称，那么，我们不把木头石头、金银当作神明，我们便是不敬虔了。但若二者相距最远——如同艺术家和他所用的材料一样遥远——为何还要追究我们呢？因为正如陶工与泥土（物质是泥土，艺术家是陶工），天主与物质的关系也是如此：天主是世界的\OriginalTerm{创造者}{Framer}，而物质为祂艺术的旨意而服务。但正如泥土若无艺术，自己不能成为器皿；同样，能接受一切形式的物质，若离了创造者天主，也不能获得区分、形状和秩序。我们不认为陶器比制造它的陶工更贵重，也不认为玻璃金器比制作它们的人更贵重；若器皿因艺术而优雅，我们赞美的乃是工匠，器皿的光荣也归于他；物质与天主也是如此——世界井然有序的光荣与尊荣，按正义不属于物质，而属于天主，即物质的\OriginalTerm{创造者}{Framer}。因此，若我们将物质的各种形式当作神明，我们便似乎是毫无对真天主的感知，因为我们竟将可分解、可朽坏之物与永恒的事物等同起来。\par

% structure: p0042 source=h2 target=section reason=relative-heading-rank; base=h2

\section{第十六章 基督徒并不敬拜宇宙。}

世界无疑是美丽的，无论在大小上还是在各部分的安排上，无论是斜圈上的还是北方的，以及它的球形形式，都出类拔萃。然而，我们应当敬拜的并非世界，而是它的\OriginalTerm{创造者}{Artificer}。因为当你们的臣民来到你们面前时，他们不会忽略向你们这些将供给他们所需的一切的统治者和主人致敬，却注目于你们宫殿的宏伟；但如果他们偶然来到王宫，他们会对它美丽的结构投以赞赏的一瞥：但他们向你们自己表示尊敬，视你们为\Quote{万有中的万有}。你们这些君主的确为自己建造和装饰宫殿；但世界被创造并非因为天主需要它；因为天主自己是祂自己的一切——不可接近的光，完美的世界，灵、能力、理性。因此，如果世界是音律和谐的乐器，且按精确的节拍运动，我所敬拜的便是那赋予和谐、弹奏音符、唱出协调旋律的\OriginalTerm{存在者}{Being}，而非乐器本身。因为在音乐比赛中，裁判不会绕过七弦琴演奏者而给七弦琴加冕。那么，无论世界如\Person{柏拉图}所言，是神圣艺术的产物，我欣赏它的美，而敬拜它的\OriginalTerm{创造者}{Artificer}；还是如\OriginalTerm{逍遥学派}{Peripatetics}所断言，世界是天主的本质和身体，我们也不忽略敬拜那位身体运动的根源——天主，并下降归向\Quote{那贫乏软弱的元素}，去敬拜那（按他们的说法）\OriginalTerm{不可感受的空气}{impassible air}，\OriginalTerm{有感受性的物质}{passible matter}；或者，若有人认为世界的各个部分是天主的大能，我们不接近去敬拜那些大能，而是敬拜它们的创造者与主。我不向物质祈求它所不能给予的，也不越过天主去敬拜那些元素，因为它们只能做它们所受命的；因为纵使它们因着其创造者的艺术而看起来美丽，然而它们仍具有物质的本性。对此，\Person{柏拉图}也作见证：\Quote{因为，}他说，\Quote{那被称为天和地的，从父那里领受了众多福分，但仍分有形体；因此，它不可能免于变化。}因此，如果我因艺术而赞赏诸天与元素，却不把它们当作神明来敬拜，因为我知道朽坏的定律临在它们身上，那么我如何能称那些我所知其创造者为人的对象为神明呢？我恳请你们留意听我对此略说几点。\par

% structure: p0044 source=h2 target=section reason=relative-heading-rank; base=h2

\section{第十七章 诸神的名字及其偶像不过是新近的。}

一位辩护者必须提出比我迄今所给出的更为精确的论据，既要论及诸神的名字，以证明它们起源晚近，也要论及它们的形象，以证明它们可以说是仅仅出现于昨日。然而，你们自己对这些事十分熟悉，因为你们精通各门知识，且比任何人都更熟谙古代事物。因此我断言，正是\Person{俄耳甫斯}、\Person{荷马}和\Person{赫西俄德}给那些他们称为神的存在赋予了谱系和名字。\Person{希罗多德}也如此见证。\Quote{他说道，\InnerQuote{我的意见是，赫西俄德和荷马比我早了四百年，仅此而已；正是他们为希腊人编造了神谱，给诸神起了名字，分配了各自的荣耀和职能，并描述了他们的形态。}} 再者，当雕塑、绘画和雕刻尚不为人所知时，神的形象根本不被使用；直到\Person{萨米亚斯}（萨摩斯人）、\Person{克拉托}（西锡安人）、\Person{克莱安特斯}（科林斯人）和那位科林斯少女出现，形象才流行起来。那时，萨米亚斯发明了轮廓画，他在阳光下勾画了一匹马；克拉托发明了绘画，他在一块涂白的板子上用油彩勾画了一男一女的轮廓；而塑造浮雕形象的技艺（\OriginalTerm{浮雕像制作}{\GreekText{κοροπλαθική}}）则由那位少女发明，她爱上了一个人，在那人熟睡时，将其影子描在墙上，她的父亲（一位陶工）欣喜于轮廓的逼真，便将草图雕刻出来并填上黏土：这尊像至今仍保存在\Place{科林斯}。此后，\Person{代达罗斯}和\Person{狄奥多鲁斯}（米利都人）进而发明了雕刻和雕塑艺术。因此你们可以明白，形象再现和制作神像的时代开始得如此之晚，以至于我们能够说出每一位特定神的工匠之名。例如，\Place{以弗所}的\Person{阿耳忒弥斯}像，以及\Person{雅典娜}像（或更确切地说，\Person{阿特拉}像，因为那些更习惯于秘仪风格的人这样称呼她；用橄榄木制成的古老神像也以此命名），还有这位女神的坐像，均为\Person{代达罗斯}的门徒\Person{恩多欧斯}所造；\Person{皮提亚}神像是\Person{狄奥多鲁斯}和\Person{特勒克利斯}的作品；\Person{得洛斯}神和\Person{阿耳忒弥斯}像出自\Person{特克泰俄斯}和\Person{安格利奥}之手；\Place{萨摩斯}和\Place{阿尔戈斯}的\Person{赫拉}像出自\Person{斯米利斯}之手，其他雕像则由\Person{菲狄亚斯}制作；\Place{克尼多斯}的妓女\Person{阿佛洛狄忒}像是\Person{普拉克西特列斯}的作品；\Place{埃皮达鲁斯}的\Person{阿斯克勒庇俄斯}像则是\Person{菲狄亚斯}的手艺。总之，没有一尊像可以说不是人手所造。那么，如果这些是神，它们为何不从起初就存在呢？事实上，它们为何比制造它们的人更年轻？为何它们需要人和技艺的帮助才能存在呢？它们不过是泥土、石头、物质和巧妙的技艺罢了。\par

% structure: p0046 source=h2 target=section reason=relative-heading-rank; base=h2

\section{第18章 如同诗人们所承认的，诸神自身也是被造的。}

然而，由于有些人断言，尽管这些只是形象，却存在它们为之而造的诸神；并且，呈献于这些形象前的祈求和献祭，乃是指向并实为献给诸神的；而且没有其他途径可以接近他们，因为\par

\Quote{凡人欲面见\newline{}一位可见的神，实属艰难；}\par

再者，他们为了证明事实确是如此，便举出某些像拥有的异能，那么让我们来考察依附于其名字的能力。最伟大的皇帝们啊，在我开始讨论之前，我恳请你们宽容我，容我提出真实的见解；因为我无意揭示偶像的虚妄，而是要驳斥对我们散播的诽谤，从而为我们的生活方式提供理由。愿你们通过自我省察，也能发现天上的国度！因为正如万有都臣服于你们——从天上承受国度的父与子——（\BibleQuote{君王的心在\Person{天主}手中}，\BibleRef{箴 21:1} \Person{圣神}说），照样，万有也同样臣服于独一的\Person{天主}和从祂而出的\OriginalTerm{圣言}{Logos}，即\Person{圣子}，我们体认\Person{圣子}与祂不可分离。我特别恳请你们仔细思想这一点。他们声称，诸神并非从起初就有，每一位都像我们一样进入存在。在这一点上，他们全体一致。\Person{荷马}说到\par

\Quote{古老的\Person{俄刻阿诺斯}，\newline{}诸神之父，和\Person{泰西丝}；}\par

以及\Person{俄耳甫斯}（不仅如此，他首先创造了他们的名字，叙述了他们的诞生，讲说了每一位的事迹，且他们相信他比其他人更真实地处理神圣事物，\Person{荷马}本人在多数事上，尤其在涉及诸神方面，也效法他）——他也将他们的最初起源定为来自水：——\par

\Quote{\Person{俄刻阿诺斯}，万物的起源。}\par

因为据他说，水是万物的源头，从水形成了泥，从两者生出一动物，一条龙，其上长有狮子的头，两头之间有一神的脸，名叫\Person{赫拉克勒斯}和\Person{克洛诺斯}。这位\Person{赫拉克勒斯}生出一个巨大无比的蛋，蛋满盈后，因生它的神的强力摩擦而裂为两半，上半部分成为\OriginalTerm{天}{\GreekText{οὐρανός}}，下半部分成为\OriginalTerm{地}{\GreekText{γῆ}}。女神\Person{盖亚}更以身体出现；\Person{乌拉诺斯}与\Person{盖亚}结合，生了女性：\Person{克罗托}、\Person{拉刻西斯}、\Person{阿特罗波斯}；以及男性：百臂的\Person{科提斯}、\Person{古革斯}、\Person{布里阿柔斯}，和独眼巨人\Person{布戎忒斯}、\Person{斯忒洛珀斯}、\Person{阿尔格斯}；他得知自己将被子女逐出权位，便捆绑了他们并投入\Place{塔尔塔罗斯}；\Person{盖亚}因此大怒，就生出了\Person{泰坦众神}。\par

\Quote{如神的\Person{盖亚}为\Person{乌拉诺斯}生下\newline{}众子，以\Person{泰坦}之名著称，\newline{}因他们向\Person{乌拉诺斯}施行报复，\newline{}那威严、闪烁着星冠的。}\par

% structure: p0055 source=h2 target=section reason=relative-heading-rank; base=h2

\section{第19章 哲学家们在关于诸神的事上与诗人们意见一致。}

他们的神祇与宇宙万有的存在，其开端便是如此。我们对此又当作何解释呢？因为凡被赋予神性的事物，都被认为是从起初就存在的。因为，如果它们曾是生成而来，先前并无存在，正如那些论述神明的人所言，它们便不是真实的存在。因为，一物若非非受造且永恒，便是受造且可朽坏的。我的看法与哲学家们并无二致。\Quote{何者是那永恒常在、无有起源的？何者又是那被生成、却又从未真实存在的？}柏拉图在论及\OriginalTerm{可理知者}{intelligible}与\OriginalTerm{可感者}{sensible}时教导说：那永恒常在的，即\OriginalTerm{可理知者}{intelligible}，是非生成的；而那非真实存在的，即\OriginalTerm{可感者}{sensible}，则是生成的，有始有终。同样，\OriginalTerm{斯多葛学派}{Stoics}也说，万物将被焚毁并再次存在，世界将获得另一个开端。但是，如果按照他们的说法，存在着双重原因：一个主动且统御的，即\OriginalTerm{天主眷顾}{providence}；另一个被动且变化的，即\OriginalTerm{质料}{matter}；然而，即便世界处于\OriginalTerm{天主眷顾}{Providence}的照料之下，因为它既为受造，便不可能保持同一状态——那么，这些并非自有、而是被生成的神明，其构造又如何能永存呢？既然这些神明的构造来源于\OriginalTerm{水}{water}，他们又比\OriginalTerm{质料}{matter}优越在哪里呢？然而，按照他们的说法，就连水也不是万物的本源。从单一而同质的元素中，又能构成什么呢？况且，质料需要一位\OriginalTerm{工匠}{artificer}，工匠也需要质料。因为没有质料或工匠，怎能造出形相呢？同样，若说质料先于天主而存在，也是不合理的；因为\OriginalTerm{动力因}{efficient cause}必然先于其所造之物而存在。\par

% structure: p0057 source=h2 target=section reason=relative-heading-rank; base=h2

\section{第二十章 对神明的荒谬描绘}

假若他们神学的荒谬之处仅限于声称神明是受造的，且其构造源自水，那么，鉴于我已证明凡被造之物无不有朽坏之虞，我自可继续应对余下的指控。但是，一方面，他们描绘了这些神明的形体：例如，将\Person{赫拉克勒斯}说成一条盘曲之龙形的神明；说其他神明有百臂之手；又如\Person{宙斯}与其母\Person{瑞亚}所生的女儿；或如\Person{得墨忒耳}，生来有两只正常的眼睛，额上另生两只，颈后长有动物的面孔，还生有双角，以致\Person{瑞亚}被这怪婴吓到，弃之而逃，没有给她\OriginalTerm{乳头}{\GreekText{θηλή}}，由此她秘仪中被称为\Person{阿忒拉}，而通常则称\Person{珀耳塞福涅}和\Person{科瑞}，尽管她与\Person{雅典娜}并非同一，后者因眼瞳而被称为\Person{科瑞}——另一方面，他们又描述了他们自认为可敬的丰功伟绩：例如，\Person{克洛诺斯}如何阉割其父，并掷之于战车之下；如何杀害自己的子女，吞食其中的男婴；以及\Person{宙斯}如何捆绑其父，抛入\Place{塔尔塔罗斯}，正如\Person{乌拉诺斯}对其子嗣所为，并与\OriginalTerm{提坦}{Titans}争战以夺取统治权；又如何因母\Person{瑞亚}拒绝与其婚配而逼迫她，她化为雌龙，他自己则变为雄龙，用所谓\OriginalTerm{赫拉克勒斯之结}{Herculean knot}将她捆绑，遂其心意，其象征便是\OriginalTerm{赫耳墨斯之杖}{rod of Hermes}；再者，他如何奸污其女\Person{珀耳塞福涅}，此次亦化为龙形，并成为\Person{狄俄尼索斯}之父。面对这般叙述，我至少必须这样说：在这种记载中，有何可取或有用之处，竟使我们必须相信\Person{克洛诺斯}、\Person{宙斯}、\Person{科瑞}等为神明呢？是那些关于他们形体的描绘吗？试问，哪个有判断力和思辨力的人，会相信一条蝰蛇是由一位神所生（诚如\Person{俄耳甫斯}所言：\par

\Quote{然而自那神圣的胎中，\Person{法涅斯}生出\newline{}另一后嗣，狰狞而凶猛，\newline{}望之乃可怖的蝰蛇，其首\newline{}生有毛发：面容姣好；其余部分，\newline{}自颈以下，则现出可畏\newline{}之龙形}）；\par

或者，谁又会承认，\Person{法涅斯}本人，这位首生之神（因为他是从卵中生出的），竟拥有龙的身躯或形状，或被\Person{宙斯}吞下，以便宙斯变得巨大无匹？因为若他们与最低贱的走兽毫无区别（显而易见，\OriginalTerm{神性}{Deity}必然异于属地之物及源于质料之物），他们便不是神。那么，请问，我们怎能以祈求者之身去亲近他们，若其起源与牲畜相仿，自身具走兽之形，且丑陋不堪入目？\par

% structure: p0061 source=h2 target=section reason=relative-heading-rank; base=h2

\section{第二十一章 归与神明的邪淫爱欲}

但若有人辩称，他们仅具血肉之形，拥有血液与种子，以及愤怒和性欲之情，即便如此，我们仍须视此类断言为荒谬可笑；因为神明中既无愤怒，亦无欲望和嗜好，更无繁衍之子种。那就让他们具血肉之形吧，但让他们超越暴怒与忿恨，免得\Person{雅典娜}被看见\par

\Quote{怒火中烧，对\Person{朱庇特}心怀愤恨；}\par

也不让\Person{赫拉}如此显现：——\par

\Quote{\Person{朱诺}的胸中\newline{}盛不下她的狂怒。}\par

让他们超越忧伤：——\par

\Quote{我双眼目睹哀恸之景：我所爱之人\newline{}绕城奔逃！我的心\newline{}为\Person{赫克托耳}而悲痛。}\par

因为即便人类中，凡纵情于愤怒与忧伤者，我亦称其为粗鲁愚拙。然而当那\Quote{人与诸神之父}为他的儿子哀悼时——\par

\Quote{哀哉，哀哉！命运竟定意，我最亲爱的\newline{}\Person{萨耳珀冬}，须倒在\Person{帕特洛克罗斯}手下；}\par

且在其哀悼之时，竟无能救他脱离险境：——\par

\Quote{朱庇特之子，朱庇特却未护佑他；}\par

谁不谴责那些人的愚昧呢？他们迷恋此类传说，以神明为爱慕，倒不如说是活在没有真神之中。让他们具血肉之形吧，但莫让\Person{阿佛洛狄忒}被\Person{狄俄墨得斯}伤及身体：——\par

\Quote{\Person{堤丢斯}骄傲的儿子，\Person{狄俄墨得斯}，\newline{}伤了我；}\par

也莫让她的心灵被\Person{阿瑞斯}所伤：——\par

\Quote{我，笨拙的我，她轻蔑不顾；却将妩媚\newline{}献与那俊美的淫棍，威武的战神}\newline{}\Quote{兵器刺穿了肌肉。}\par

那在战场上可畏惧者，\Person{宙斯}对抗\OriginalTerm{提坦}{Titans}时的盟友，竟显得比\Person{狄俄墨得斯}更软弱：——\par

\Quote{他狂怒，如\Person{玛尔斯}挥舞长矛。}\par

嘘！\Person{Homer}，神从不发怒。但你却将这位神描绘成沾染鲜血、是凡人的祸害：——\par

\Quote{\Person{Mars}，\Person{Mars}，凡人的祸害，沾满鲜血；}\par

你又说他的奸淫和他的捆锁：——\par

\Quote{于是，他毫不迟疑，领着那迷醉的美人，\newline{}心荡神驰，沉入知情的床榻。\newline{}罗网骤然落下。}\par

关于诸神，他们岂不是大量倾泻这类不虔敬的东西吗？\Person{Ouranos}被阉割；\Person{Kronos}被捆绑，并被推入\Place{Tartarus}；\Person{Titans}反叛；\Person{Styx}战死：是的，他们甚至将诸神描绘成有死的；他们彼此相爱；他们与人类相爱：——\par

\Quote{\Person{Æneas}，在\Place{Ida}山突出的群峰间，\newline{}不朽的\Person{Venus}为\Person{Anchises}所育。}\par

他们不是恋爱吗？他们不是受苦吗？不，真的，他们是神，欲望不能触及他们！即使一个神为执行一项神圣计划而\OriginalTerm{取肉身}{assume flesh}，他因此就成了欲望的奴隶。\par

\Quote{因为从未有过如此泛滥的爱情，\newline{}为女神或凡人，充满我的灵魂；\newline{}不为\Person{Ixion}的美妻，她生了\newline{}\Person{Pirithöus}，如诸神般睿智的谋士；\newline{}也不为整洁脚踝的少女\Person{Danäe}，\newline{}\Person{Acrisius}之女，她生了\Person{Perséus}，\newline{}众人所瞩目；也不为高贵的\Person{Phœnix}之女；\newline{}……也不为\Person{Semele}；\newline{}也不为美丽的\Person{Alcmena}；……\newline{}不，也不为金发的女王\Person{Ceres}；\newline{}也不为光明的\Person{Latona}；也不为你自己。}\par

他是受造的，是会朽坏的，在他身上毫无神的踪迹。不仅如此，他们甚至成了人的雇工：——\par

\Quote{\Person{Admetus}的宫殿，在那里我虽为神，\newline{}却忍受赞美奴仆的餐桌。}\par

他们还放牧牛群：——\par

\Quote{来到这片土地，我为他放牧牛群，\newline{}他款待我，我守护他的家。}\par

因此，\Person{Admetus}胜过了那位神。你，先知和智慧者，能为别人预见未来的事，却没有预知你所爱之人的被杀，甚至亲手杀了他，尽管他如此亲爱：——\par

\Quote{我原相信\Person{Apollo}神圣的口\newline{}充满真理，以及先知的艺术。}\par

\EditorialNote{埃斯库罗斯指责阿波罗是假先知：}\par

\Quote{正是在宴席上歌唱的那一位，\newline{}说这些话的那一位，哀哉！正是他\newline{}杀了我的儿子。}\par

% structure: p0094 source=h2 target=section reason=relative-heading-rank; base=h2

\section{第二十二章 伪托的寓意解释。}

但或许这些只是诗意的幻想，对它们有某种自然解释，比如\Person{Empedocles}的这个解释：——\par

\Quote{让\Person{Jove}为火，\Person{Juno}为生命之源，\newline{}连同\Person{Pluto}和\Person{Nêstis}，她用泪水浸洗\newline{}人类的泉源。}\par

因此，如果\Person{Zeus}是火，\Person{Hera}是土，\Person{Aïdoneus}是气，\Person{Nê stis}是水，而这些是元素——火、水、土、气——那么它们中没有一个是神，\Person{Zeus}不是，\Person{Hera}不是，\Person{Aïdoneus}也不是；因为它们的组合和起源来自于被\OriginalTerm{天主}{God}分开的物质：——\par

\Quote{火、水、土，以及气的轻柔高度，\newline{}以及与这些的和谐。}\par

这些事物若无和谐便不能存留，在纷争中便会毁灭：既然如此，怎可称它们为神呢？\Person{恩培多克勒}认为，友爱具有治理的才能，复合之物被治理，而善于治理的拥有统御权；因此，倘若我们将受治者与治理者的能力等同起来，便会在不知不觉中将可朽坏、变动不居、变幻无常的物质，置于与非受造、永恒且永远自相一致的天主同等的地位。据斯多葛派所言，\Person{宙斯}是自然中炽热的部分；\Person{赫拉}是\OriginalTerm{空气}{\GreekText{ἀήρ}}——若将这个名字自身相连，即标示此意；\Person{波塞冬}是饮料（水，\OriginalTerm{饮料}{\GreekText{πόσις}}）。然而，这些事物被不同的人以不同的方式解释为自然对象。有人称\Person{宙斯}为双性的阳气；另有人说他是带来温和天气的季节，正因如此，他才独自逃脱了\Person{克罗诺斯}的吞噬。但对斯多葛派，我们可以说：如果你们承认一位至高、非受造且永恒的天主，并认为有多少物质的变化，就有多少复合的身体，且说天主的神弥漫于物质中，随其变化而获得不同的名称，那么物质的各种形式就成了天主的身体；但当元素在宇宙大火中毁灭时，名称必然与形式一同消亡，唯有天主的神存留。那么，谁能相信那些随物质而变化、与朽坏相关联的形体是神呢？但对于那些说\Person{克罗诺斯}是时间、\Person{瑞亚}是大地、她因\Person{克罗诺斯}而怀孕生养，因此被视为万物之母；说他生养并吞噬自己的子女；说阉割是男女的交合，截断种子投入子宫，生出具有性欲的人，而这性欲就是\Person{阿佛洛狄忒}；说\Person{克罗诺斯}的疯狂是季节的变换，毁灭一切有生无生之物；说锁链和\Place{塔尔塔罗斯}是时间，随季节改变而消逝的人，我们回答说：若\Person{克罗诺斯}是时间，他便有变化；若是季节，他便有轮转；若是黑暗、霜冻，或自然的潮湿部分，这些都非恒常；但天主是不死不灭、不动不变的，因此，\Person{克罗诺斯}及其象征都不是天主。至于\Person{宙斯}：若他是气，由\Person{克罗诺斯}所生，其中阳性部分称为\Person{宙斯}，阴性称为\Person{赫拉}（因此她既是姐妹又是妻子），他便受制于变化；若是季节，便有轮转：但天主既不改变，也不转移。但我何必再絮聒不休，徒耗你们的耐心呢？你们熟知那些将这些事物化解为自然要素的人各自说了什么，也熟知不同的作者对自然的看法，以及他们对\Person{雅典娜}的说法——他们断言\Person{雅典娜}是弥漫于万物中的\OriginalTerm{智慧}{\GreekText{φρόνησις}}；对\Person{伊西斯}的说法——他们称她为一切时间的诞生（\OriginalTerm{一切时间的自然}{\GreekText{φύσις} \GreekText{αἰῶνος}}），万物由她而生，藉她而存；对\Person{奥西里斯}的说法——他被兄弟\Person{堤丰}杀害后，\Person{伊西斯}携子\Person{奥鲁斯}寻找他的肢体，找到后以坟墓尊崇，这坟墓至今仍被称为\Person{奥西里斯}的墓。因为他们徘徊于物质的形式之间，错过了那只能由理智观照的天主，反倒将元素及其各个部分神化，在不同的时间赐予它们不同的名称：例如，将谷物的播种称为\Person{奥西里斯}（因此他们说，在秘仪中，当他们寻见他身体的肢体或果实，便对\Person{伊西斯}这样祝祷：\Quote{我们找到了，我们向你祝贺}），将葡萄的果实称为\Person{狄俄尼索斯}，葡萄树本身称为\Person{塞墨勒}，太阳的热力称为雷电。然而，事实上，那些将这些神话归于实际神祇的人，非但没有增添其神性，反而适得其反；因为他们未曾察觉，正是他们为这些神所做的辩护，确认了那些关于他们的传说。\Person{欧罗巴}、公牛、天鹅以及\Person{勒达}，与大地和空气有何相干，竟能将\Person{宙斯}与之的可憎交合解作大地与空气的交合？他们未能发现天主的伟大，又不能凭理智上升高天（因为他们与天上的境界无缘），遂在物质的形式中憔悴消耗，根植于大地，将元素的变化神化：正如同有人将他所乘的船置于舵手的位置。然而，正如船虽配备齐全，若无舵手便毫无用处，元素虽排列整齐，若无天主的眷顾也毫无功用。船不会自行航行；元素若无创造者，也不会运动。\par

% structure: p0100 source=h2 target=section reason=relative-heading-rank; base=h2

\section{泰勒斯与柏拉图的意见。}

然而，你或许会说，既然你的理解力超越所有人，那么，若那些我们为之立像的并非神，为何有些偶像竟能显出大能呢？因为无生命的、不能运动的形象，若无推动者，是不能自行做任何事的。我们绝不否认，在各地、各城、各族中，确有某些功效借着偶像的名发生。然而，即便有人因此受益，另有人反遭损害，我们也绝不可认为那产生这些效果的便是神。但我已仔细探究：为何你以为偶像有此大能？又是谁盗用其名产生这些效果？为了说明产生偶像效果的是谁，且他们并非神，我有必要引证一些哲学家的见解为佐证。首先，\Person{泰勒斯}，据那些准确研究过他学说的人报告，将[高级存在]分为：天主、魔鬼和英雄。他认天主为世界的\OriginalTerm{理智}{\GreekText{νοῦς}}；以魔鬼指\OriginalTerm{有灵魂的存在}{\GreekText{ψυχικαί}}；以英雄指人分离后的灵魂，好的即是善的灵魂，恶的便是无价值的。\Person{柏拉图}，在其他方面虽未予赞同，也同样将[高级存在]分为：非受造的天主、由那非受造者所生、为装饰诸天、行星及恒星而造的存在，以及魔鬼；关于这些魔鬼，他认为不宜亲自谈论，却以为当听从那些曾讲述他们的人。\Quote{要谈论其余的魔鬼，并知晓其起源，是超越我们能力的；但我们应当相信那些曾讲述过的人，即众神的子孙，如他们所言——他们必定深知自己的祖先：因此，即便他们说起话来没有可信或确切的证据，也不可不信神的儿女们；但他们既声称讲述自家事务，我们遵循习俗，便理应相信他们。既然如此，我们就当像他们那样，持守并谈论众神本身的起源。从\Person{盖亚}和\Person{乌拉诺斯}生出了\Person{俄刻阿诺斯}和\Person{忒堤斯}；从这些又生出\Person{福耳库斯}、\Person{克洛诺斯}、\Person{瑞亚}等；从\Person{克洛诺斯}和\Person{瑞亚}则生出了\Person{宙斯}、\Person{赫拉}及其他一切，据我们所知，他们全都互称兄弟；此外还有他们的其他后裔。}那么，这位曾默观永恒理智和那可凭理性领悟的天主，并宣告了祂的属性——祂的真实存在、祂本质的单纯、从祂流出的善即是真理，且论述了原初大能，以及\Quote{万物皆环绕万有之王，万物皆因祂而存在，祂是万物之因；}又论及二与三，说祂是\Quote{第二者运行于第二领域，第三者运行于第三领域；}——难道他会认为，探究那些据称由可感之物（即地与天）所生者的真相，竟是超越其能力的吗？这是绝不可信的。而是因为，他认为不可能相信众神会生育和被生，因为一切开始存在的，必有终结随之而来，并且（这更为困难）要改变那些不假思索便接受这些神话的大众的意见，正是因此，他才宣称自己无力知晓并谈论其他魔鬼的起源，因为他既不能承认，也无法教导神是被生的。至于他那句话：\Quote{天上的伟大君王\Person{宙斯}，驾着羽翼战车，行进在前，安排统御万有，众神与魔鬼的大军紧随其后。}这并不是指那位据称由\Person{克洛诺斯}所生的\Person{宙斯}；因为在此，这名是赋予宇宙的创造者的。\Person{柏拉图}自己表明了这点：因无法用其他适宜的名号来称呼祂，便借用了通俗的名字，并非将其当作天主的专名，而是为了清楚表达，因为人不可能向所有人完满地讲述天主；同时他加上了\Quote{伟大}这个修饰，好将属天的与属地的、非受造的与受造的区别开来，后者比天地还要年轻，且比那些将他偷走、以免他被父亲杀死的\Person{克里特人}更年轻。\par

% structure: p0102 source=h2 target=section reason=relative-heading-rank; base=h2

\section{第24章 论天使与巨人}

在向你们这些已探究一切知识领域的人说话时，何需提及诗人，或详察其他种类的见解呢？只消如此说便够了。倘若诗人和哲人们不承认有一位天主，且关于这些神明，他们的意见并非有些认为是魔鬼，有些认为是物质，有些认为曾是人——那么，我们受逼迫或许还有几分理由，因我们使用一种区分天主与物质及二者本性的言论。因为，我们承认一位天主，一位圣子——祂的\OriginalTerm{圣言}{Logos}，以及一位圣神，在\OriginalTerm{本质}{essence}内合一——圣父、圣子、圣神，因为圣子是圣父的\OriginalTerm{理智}{Intelligence}、\OriginalTerm{理性}{Reason}、\OriginalTerm{智慧}{Wisdom}，而圣神是一种\OriginalTerm{流露}{effluence}，如同来自火的光；同样，我们也领会还有其他异能存在，它们统辖物质，并藉由物质行事，其中有一个特别与天主为敌：并非真有什么能与天主对立，如同恩培多克勒所说的斗争与友谊，或如星辰之显现与消逝中夜晚与白昼的对立（因为即便有什么\Emph{曾}与天主对立，它也会不复存在，其结构被天主的能力与大能摧毁），而是与天主内的\OriginalTerm{善}{good}——这善必然属于祂，并与祂共存，如同颜色与身体，没有它身体便不存在（并非作为它的一部分，而是作为一种伴随的属性与之共存，结合交融，正如同火自然为黄色，以太自然为深蓝色）——与天主内的善，我说，那环绕\OriginalTerm{物质}{matter}的灵，乃是由天主所造，正如其他天使由祂所造，并被托付管理物质及其形式，是与之为敌的。因为这正是天使的职分——为天主向祂所造并安排的事物施行照顾；如此，天主对整个宇宙享有普遍而周详的照顾，而个别部分则由受命的天使负责照料。正如同在人中间，人对于德行与恶行都有自由选择（因为除非恶与德皆在各自权下，你们既不会尊敬善人，也不会惩罚恶人；且有些人对于你们所托付的事殷勤，另有些人则不忠），在天使中也是如此。你们将观察到，有些天使，作为自由的代理者，一如受造于天主的模样，持守在天主创造并规定他们的那些事中；但有些则冒犯了自己本性的构造以及所受托的管理：就是说，这位\OriginalTerm{物质之主宰}{prince of matter}及其各种形式，以及其他置于这第一个穹苍周围的天使（你们知道，我们所说无不有见证，乃陈述先知们所宣告的事）；他们坠入对贞女不洁的爱恋，被肉情制服，他在所托之事的管理上变得疏忽而邪恶。于是，从这些恋慕贞女者，便生出了所谓的巨人。若是诗人也曾提及巨人，对此不必惊讶：世俗的智慧与属神的智慧，二者的差别正如真实与似是而非：一个属乎上天，另一个属乎大地；而按照\OriginalTerm{物质之主宰}{prince of matter}的话说——\par

\Quote{我们知道，我们常说出看似真理的谎言。}\par

% structure: p0105 source=h2 target=section reason=relative-heading-rank; base=h2

\section{诗人和哲人们曾否认\OriginalTerm{天主眷顾}{Divine Providence}。}

因此，这些自天坠落的天使，游荡于空中和地上，再不能升达天上的事理，还有巨人们的灵魂，即那些徘徊世间的魔鬼，二者行径相似，前者（即魔鬼）依其禀受之性而行，后者（即天使）则随其放纵之欲而为。然而，\OriginalTerm{物质之主宰}{prince of matter}，仅从所发生的事便可看出，施行与天主内的善背道而驰的控制与管理——\par

\Quote{时常这忧思萦绕我心，\newline{}究竟运数还是神祇主宰\newline{}人间微渺之事；不顾希望\newline{}与公义，将某些人剥夺一切\newline{}生计，驱于流放，而其他人仍\newline{}继续安享荣华。}\par

荣华与灾祸，与希望和公义相悖，使欧里庇得斯无从断定世间事务的管理属于谁，其情状令人不禁言道：——\par

\Quote{那么，目睹此景，我们怎能说\newline{}有一族神明，或服从法律？}\par

同一事由使亚里士多德宣称，天以下的事物不受\OriginalTerm{天主眷顾}{Providence}的照管，虽然天主的永恒眷顾同样关切我们这些下界之人——\par

\Quote{大地，无论意愿动与不动，\newline{}必出产草木，从而养活我的羊群，}\par

它按真理而非意见，个别地向配得者发言；所有其余事物，依照大自然的普遍构造，皆由理性之律安排。然而，因由恶神而出的恶魔之运作与行动产生了这些混乱的发作，并且进一步按物质倾向一方面、对神圣事物的亲和力另一方面，从内及外地感动人，有的这样，有的那样，作为个人与民族，各自与共同地——有些并非无名之辈遂以为这个宇宙的构成原无定序，被非理性的偶然所驱动。但他们不明白，那属于全宇宙构造的事物，无一陷入无序或受忽略，反之，每一件都由理性产生，因而并不逾越为其所规定的秩序；且人自身，就其受造于造他的那位而言，亦秩序井然，既因其原本本性具有一个对一切人皆共通的普遍特性，又因其身体构造不逾越加于其上的规律，且因其寿命之终结保持平等并同等地属于所有人；但依其自身特有的性格，以及掌管世间的君侯及其随从魔鬼的作用，他被推动、被驱使朝向这方或那方，纵然一切人共同拥有同一的心神原始构造。\par

% structure: p0113 source=h2 target=section reason=relative-heading-rank; base=h2

\section{第二十六章 魔鬼引诱人敬拜偶像}

那么，引人归向偶像的，正是前述的魔鬼，他们贪嗜祭牲的血，餂噬之；但悦于大众、其名被赋予那些像的诸神，原是人，这可从他们的历史得知。而那些冒其名行事的是魔鬼，又由他们行事之性质证明。因为有的致人阉割，如\Person{瑞亚}；有的致伤和杀戮，如\Person{阿耳忒弥斯}；\Person{陶洛女神}将一切异乡人处死。我略过那些用刀与骨鞭伤人者，也不试图记述所有种类的魔鬼；因为激人干犯逆性之事，并非神明之分。\par

\Quote{但当魔鬼图谋陷害人时，\newline{}他先对其心志施以伤害。}\par

然而，天主，作为全然美善者，永恒地行善。再者，行使能力者与那立像的对象并非同一，\Place{特洛阿斯}与\Place{帕里乌姆}提供了极有力的证据。前者有我们同时代人\Person{尼律利努斯}的像；\Place{帕里乌姆}则有\Person{亚历山大}和\Person{普洛透斯}的像：\Person{亚历山大}的陵墓和像至今仍矗立在广场。\Person{尼律利努斯}其余的像，则是一种公共装饰——如果一座城能被这类物件所装饰的话；但其中一尊据说发出神谕、医治病人，为此\Place{特洛阿得}人向这尊像献祭，并包金、挂花环。至于\Person{亚历山大}和\Person{普洛透斯}（后者，如你所知，在\Place{奥林匹亚}附近投身火中）的像，\Person{普洛透斯}的像据说同样发出神谕；而\Person{亚历山大}的像——\par

\Quote{可怜的\Person{帕里斯}，虽外貌如此俊美，\newline{}你却是女人的奴隶} ——\par

以公费向之献祭、举行节期，如同向能听的神一般。那么，是\Person{尼律利努斯}、\Person{普洛透斯}和\Person{亚历山大}借这些像行使力量，还是物质自身的本性如此？但物质乃是黄铜。黄铜本身能做什么呢？它本可被重新铸成不同的形制，一如\Person{赫洛多托斯}所记，\Person{阿玛西斯}处理脚盆的故事。\Person{尼律利努斯}、\Person{普洛透斯}和\Person{亚历山大}，对病人又有何益？因为那像据说现在所产生的作用，它在\Person{尼律利努斯}活着且患病时就已产生。\par

% structure: p0119 source=h2 target=section reason=relative-heading-rank; base=h2

\section{第二十七章 魔鬼的诡计}

那么，究竟是什么？首先，灵魂关于意见的非理性和幻想性的运动，不时产生出各种\OriginalTerm{幻象}{\GreekText{εἴδωλα}}：有些得自物质，有些则自己塑造孕育；这尤其发生在一个灵魂分沾了物质之神灵，且与之搅合的时候，因为它不仰望天上的事物和它们的造主，反而向下垂顾属土之事，全然注视大地，仿佛成了纯粹的肉血，而不再是纯神。于是，灵魂这些非理性和幻想性的运动，便在心智中产生虚妄的异象，使其疯狂地执迷于偶像。再者，当一个柔嫩易感的灵魂，对更健全的教义既无知识亦无经验，不惯于观照真理，也不惯于深思圣父和万有的造主，却又被关于自身的错误意见所熏染，那么，那些盘桓于物质、贪嗜祭牲的气味和牺牲的血、且随时准备引人入误途的魔鬼，便利用群众灵魂的这些虚幻运动；并且，占据他们的思想，使之从偶像和雕像仿佛产生虚空的异象，涌入心智；而当灵魂本身，因是不死不灭的，按理性活动，或预言未来，或医治现症，魔鬼却将这荣光归于自己。\par

% structure: p0121 source=h2 target=section reason=relative-heading-rank; base=h2

\section{第二十八章 异教神祇不过是凡人}

但或许有必要，根据已经提出的，略谈一下他们的名字。于是，\Person{希罗多德}，以及\Person{腓力}之子\Person{亚历山大}，在他给母亲的信中（据说他们两人都曾在\Place{赫利奥波利斯}、\Place{孟斐斯}和\Place{底比斯}与祭司交谈过），确认他们从祭司那里得知，诸神曾是人。\Person{希罗多德}如此说：\Quote{他们说，这些形象所代表的存在，其性质如此，它们远远算不上是神。然而，在它们之前的时代则不然；那时，埃及有诸神作统治者，与世人一同住在地上，其中一位总是凌驾于其他之上。最后一位是\Person{奥西里斯}之子\Person{何鲁斯}，被希腊人称为\Person{阿波罗}。他废黜了\Person{提丰}，作为最后一位神王统治了埃及。\Person{奥西里斯}被希腊人称为\Person{狄俄尼索斯}（\Person{巴库斯}）。} \Quote{几乎所有神的名字都是从埃及传入希腊的。} \Person{阿波罗}是\Person{狄俄尼索斯}和\Person{伊西斯}之子，正如\Person{希罗多德}同样确认的：\Quote{据埃及人说，\Person{阿波罗}和\Person{狄安娜}是\Person{巴库斯}和\Person{伊西斯}的孩子；而\Person{拉托娜}是他们的保姆和保存者。} 这些具有天上起源的存在，他们将其立为首批国王：部分因为不知晓对天主的真正崇拜，部分出于对他们统治的感激，他们便连同他们的妻子一起尊崇为神。\Quote{所有埃及人都用洁净的公牛和公犊献祭；但母牛他们不允许献祭，因为它们是\Person{伊西斯}的圣物。这位女神的雕像呈女人形，却像母牛一样有角，类似于希腊人对\Person{伊俄}的描绘。} 还有谁比那些代代相传，父传子不仅继承了祭司职分，也继承了历史的人，在作此陈述时更值得信赖呢？因为祭司的职责是使人们尊崇偶像，他们不大可能虚假地声称这些偶像曾是人。如果只有\Person{希罗多德}说，埃及人在他们的历史中把诸神当作人来谈论，正如他所说的：\Quote{关于他们的宗教，他们告诉我的，我无意复述，只提一下他们神祇的名字，那些都是微不足道的事情}，那么连\Person{希罗多德}也可能被视为编造神话的人，我们不必相信。但既然\Person{亚历山大}，以及\Person{赫耳墨斯}，别号\OriginalTerm{特里斯墨吉斯忒斯}{Trismegistus}，与他们同具永恒属性，还有无数其他的人（不用一一指名），都表明同样的事，那么诸王曾被尊崇为神，便毫无疑问了。最博学的埃及人也见证他们曾是人，这些埃及人虽说以太、大地、太阳、月亮是神，但认为其余的都是凡夫俗子，而庙宇不过是他们的坟墓。\Person{阿波罗多洛斯}也在他关于诸神的论著中主张同样的事。然而\Person{希罗多德}甚至把他们的苦难称为奥秘。\Quote{在\Place{布西里斯}城举行的\Person{伊西斯}节上的仪式，前已述及。在那里，男女大众，成千上万，在献祭结束时，为了崇敬一位我因宗教顾忌不便提及其名号的神，而捶打自己。} 如果他们真是神，那他们也是不朽的；但如果人们为他们捶打自己，而且他们的苦难被称为奥秘，那么他们就是人，正如\Person{希罗多德}自己所说：\Quote{同样，就在\Place{萨伊斯}城的这个\Person{密涅瓦}圣区内，有一个人的葬地，我认为在这样的话题中提及他是不合适的。它位于神殿后面，靠在后墙上，完全覆盖住那面墙。圣区内还有一些大石方尖碑，附近有一个湖，边缘用石头砌成。它的形状是圆形的，大小，据我看来，大约与\Place{提洛}岛上叫做\InnerQuote{环}的湖相当。就在这个湖上，埃及人夜间表演那位我避而不提其名号者的受难，他们将这种表演称为他们的奥秘。} 而且不仅展示\Person{奥西里斯}的坟墓，还展示了他的防腐处理：\Quote{当一具尸体被带给他们时，他们向送葬者展示各种用木料制成并涂上颜色以模仿自然的尸身模型。据说其中一种最为完美，是依照那位我认为在涉及此类事宜时不宜提及其名号者的样式制作的。}\par

% structure: p0123 source=h2 target=section reason=relative-heading-rank; base=h2

\section{二十九、从诗人证明同样的事。}

但在希腊人中，那些以诗歌和历史著称的人也说着同样的事情。关于\Person{赫拉克勒斯}如此说：——\par

\Quote{那无法无天的恶棍，那蛮力之人，\newline{}对上天之声充耳不闻，亵渎了社会礼仪。}\par

他的本性如此，理应发疯，理应点燃火葬柴堆将自己烧死。关于\Person{阿斯克勒庇俄斯}，\Person{赫西俄德}说：——\par

\Quote{诸神与人类强大的父亲\newline{}充满了愤怒，从\Place{奥林波斯}山巅\newline{}用烈焰霹雳掷下，击杀了\newline{}\Person{拉托娜}所钟爱之子——他对此深为恼怒。}\par

\Person{品达}也说：——\par

\Quote{但即使智慧也会被利益所陷。\newline{}手中可见的灿烂黄金之贿赂\newline{}甚至也败坏了他：因此，\Person{克罗诺斯}之子\newline{}用双手迅速扼止了他的气息，\newline{}并以一道火雷确保了他的灭亡。}\par

因此，要么他们是神，不贪恋黄金——\par

\Quote{哦，黄金，必有一死的人最美好的奖品，\newline{}无论是母亲，还是亲爱的儿女，\newline{}都无法在喜悦中与之相匹}\par

因为天主原无需任何事物，且超越肉欲，他们并未死去；或者，他们生而为人，因无知而成为邪恶的，并被贪爱金钱所胜。还有什么需要我说，或提及\Person{卡斯托尔}、\Person{波吕丢刻斯}、\Person{安菲阿拉俄斯}呢？他们可以说刚刚出生不久，是人之子，却被视为神，甚至在\Person{伊诺}经历了疯狂及其后的苦难后，他们想象她已成为女神？\par

\Quote{航海者将称她名为\Person{琉科忒亚}。}\par

至于她的儿子：——\par

\Quote{威严的\Person{帕莱蒙}，水手们将呼求。}\par

% structure: p0136 source=h2 target=section reason=relative-heading-rank; base=h2

\section{三十、将神性归于人的原因。}

可憎且为神所憎恶的人尚且享有神明之誉，而\Person{德尔刻托}的女儿\Person{塞弥拉弥斯}，一个淫荡而染血的女人，竟被尊为\Place{叙利亚}女神；并且，倘若因\Person{德尔刻托}之故，\Place{叙利亚}人敬拜鸽子和\Person{塞弥拉弥斯}（因为，一件不可能的事，一个女人变成了鸽子：这故事见于\Person{克特西亚斯}），那么，若有人因其统治与主权而被其人民称为神明，又有何奇怪呢（\Person{西比尔}，\Person{柏拉图}也曾提及她，说：——\par

\Quote{那时是第十代，\newline{}具有语言能力的人类，自从洪水\newline{}冲毁先前时代的人类以后，\newline{}\Person{克洛诺斯}、\Person{伊阿珀托斯}和\Person{提坦}为王，\newline{}人们称他们为\Person{乌拉诺斯}和\Person{盖亚}\newline{}最高贵的儿子，并如此命名，\newline{}因为在具有语言能力的凡人中\newline{}他们是最早的}）；\par

还有因其力量而获封神者，如\Person{赫拉克勒斯}与\Person{珀耳修斯}；又有因其技艺而受封者，如\Person{阿斯克勒庇俄斯}？因此，那些或被臣民赋予尊荣，或由统治者自身[攫取]尊荣者，便获得了神明之名，有的是出于畏惧，有的是出于报复。这样，\Person{安提诺乌斯}因你们祖先对臣民的仁惠而获敬拜，被奉为神。但后来的人不经省察便接受了这种敬拜。\par

\Quote{克里特人总是说谎；因为，王啊，\newline{}他们为你建造了一座坟墓，而你并未死去。}\par

\Person{卡利马科斯}啊，你虽然相信\Person{宙斯}的诞生，却不相信他的坟墓；你以为在掩盖真相，实际上你甚至向无知者宣告他已经死去；你若看见山洞，便想起\Person{瑞亚}的分娩；但当你看见棺材时，你却对他的死亡投下阴影，却不考虑唯有那非受生的\OriginalTerm{天主}{God}才是永恒的。因为，要么大众和诗人们讲述的关于众神的故事不可信，对他们的尊崇是多余的（因为那些故事若为虚假，众神便不存在）；或者，若那些诞生、艳史、谋杀、偷盗、阉割、雷电皆为真实，那么他们便不复存在，他们既已生出，便告终结，而先前本不存在。诗人们撰写这些故事原为增加敬意，我们又根据什么原则信其一部分而不信另一部分呢？因为，那些因这些故事而被视为神明，并竭力将其行为描绘为值得崇敬的人，当然不可能编造出他们的苦难。因此，我们并不是无神论者，我们承认\OriginalTerm{天主}{God}——这宇宙的创造者——及祂的\OriginalTerm{圣言}{Logos}，这一点已尽力加以证明，纵然未必与课题的重要性相称。\par

% structure: p0142 source=h2 target=section reason=relative-heading-rank; base=h2

\section{第三十一章 驳斥对基督徒的其他指控}

此外，他们还编造了关于我们举行不虔敬的宴饮和禁止的男女交合的故事，一方面为使自己的仇恨显得有理有据，另一方面认为，可以藉暴虐或恐怖使我们偏离自己的生活方式，或藉指控的重大使官长变得严酷无情。然而，对于知道自古以来——而非仅在我们这个时代——便是邪恶攻击美善的人来说，他们徒劳无功。正如\Person{毕达哥拉斯}与三百余人一同被烧死；\Person{赫拉克利特}和\Person{德谟克利特}被放逐，一人被逐出\Place{厄弗所}人的城，另一人被逐出\Place{阿布德拉}，皆因被控癫狂；\Place{雅典}人更将\Person{苏格拉底}判处死刑。但正如他们不会因大众的意见而在品德上有丝毫减损，同样，某些人不加分辨的诽谤也不会在我们生活的正直上投下任何阴影，因为我们在\OriginalTerm{天主}{God}前享有美名。不过，我仍然要回应这些指控，尽管我确信，通过已说的话，我已经在你们面前为自己洗清了。因为你们在才智上卓绝众人，你们知道，那些以\OriginalTerm{天主}{God}为准则引导其生活的人，使我们在祂面前无可指摘、无可非议，这样的人连最微小的罪过都不会起念。因为如果我们只相信今生唯一，那我们或许会因受制于血肉，或被利欲或肉欲所胜，而有犯罪之嫌疑；但既然我们知道，\OriginalTerm{天主}{God}见证我们日夜之所思所言，而祂本身就是光，洞悉我们心中的一切，我们便确信，当我们离开今世以后，我们将进入另一种生命，比今世更好，是属天的，而非属地的（因为我们将在\OriginalTerm{天主}{God}身旁，与\OriginalTerm{天主}{God}同在，灵魂远离一切变化与痛苦，不再如血肉，尽管我们仍有肉身，而是如同属天的精神体），或与其余之人一同堕落，进入更坏的生命并在火中；因为\OriginalTerm{天主}{God}并没有把我们造得像羊或负重牲口一般，仅为附带的工作，且让我们灭亡归无。基于这些理由，我们不可能愿意作恶，也不可能将自己交予那位伟大的审判者去受惩罚。\par

% structure: p0144 source=h2 target=section reason=relative-heading-rank; base=h2

\section{第三十二章 基督徒高尚的道德}

然而，他们编造有关我们的故事，就像他们讲述自己神明的生活事迹并造作奥秘那样，这实在不足为奇。但是，如果他们真想谴责无耻而淫乱的交合，那么他们要么该憎恶\Person{宙斯}——他与母亲\Person{瑞亚}、女儿\Person{科瑞}生子，又娶自己的姐妹为妻；要么该憎恶这些故事的编造者\Person{俄耳甫斯}，他让\Person{宙斯}比\Person{梯厄斯忒斯}本人更不圣洁、更可憎；因为梯厄斯忒斯为了求取神谕、夺取王位并施行报复，才玷污了自己的女儿。然而我们非但不作淫乱的交合，就连一个淫邪的目光也不容许放纵。因为祂曾说：\BibleQuote{因为凡注视妇女，有意贪恋她的，他已在心里奸淫了她。}\BibleRef{玛 5:28} 因此，那些被禁止观看超过天主造眼之目的的事物——眼睛本是为我们作光明的，若以它追求别的事物，淫邪的目光便是奸淫；那些连自己的思想都要受审问的人，谁能怀疑他们躬行节制呢？因为我们所交代的事，不依据人制定的法律，恶人尚可逃避（我在前面对你们，至尊的君主，已经证明我们的教义乃源自天主的教导），但我们有一条法律，它以对待近人如同自己作为正直的尺度。也正因此，我们按年龄分别看待：有的视如儿女，有的当做兄弟姊妹，对年长者则致以父母般的敬重。那么，对于那些我们称以兄弟、姊妹和其他亲属名号的人，我们极其小心地保持他们身体的圣洁无损；因为\OriginalTerm{圣言}{Logos}再次对我们说：\Quote{凡有人因着第二吻而感到愉悅，便是犯了罪；}接着又说：\Quote{因此，亲吻或问候，必须极之谨慎，若在其中掺杂丝毫思想上的污秽，便将我们摒于永生之外。}\par

% structure: p0146 source=h2 target=section reason=relative-heading-rank; base=h2

\section{第三十三章 论基督徒在婚姻上的贞洁}

因此，既怀着永生的希望，我们便轻视今生的事物，甚至灵魂的享乐。我们各人按着我们所立的规条，以那合法迎娶的妻子为伴侣，并且只为生育子女。如同农夫把种子撒在地里，等待收获，而不在上面过度撒种，同样，生育子女便是我们容许欲望的界限。不仅如此，你会发现我们中间有许多男女，为着希望能与天主更亲近地共融，终身不婚，守贞至老。如果守童贞和阉人的身份使人更亲近天主，而放纵肉欲与邪念却使人远离祂，那么，我们既然连邪念都躲避，就更加憎恶其行了。因为我们所致力宣扬和教授的，不是言词之学，而是行为——即人当或守其本性，或满足于一夫一妻的婚姻；因为第二次婚姻只是伪装的奸淫。因为祂曾说：\BibleQuote{凡休妻而另娶的，便是犯奸淫；}\BibleRef{玛 19:9} 祂不许人休弃已经破身的妻子，也不许再娶。因为凡离弃元配的人，即使她已经去世，也是一个伪装的奸淫者，他在违抗天主之手，因为起初，天主造了一男一女，并为种族繁衍而设立了肉与肉的最紧密结合；凡再婚者，便是拆散了这一结合。\par

% structure: p0148 source=h2 target=section reason=relative-heading-rank; base=h2

\section{第三十四章 基督徒与控告者之间天壤之别的道德}

然而，虽然我们品行如此（噢！我何必说那些不堪言说的事呢？）人们所捏造关于我们的事，正应验了那句谚语：\Quote{妓女责备贞妇。} 因为那些人开设卖淫的市场，为青年设立种种下流娱乐的淫窟——他们甚至不避开男性，男与男行那骇人听闻的可憎之事，以种种方式玷污一切高贵俊美的躯体，就这样羞辱天主美妙的造化（因为世间的美丽并非自生，而是由天主的手和旨意差遣而来）——这些人，我要说，将自己明知的事来毁谤我们，并将之归于他们的神明，夸耀这些事为高贵的行为，配得上神明。这些奸淫者和恋童者，却诋毁阉人和一次婚姻者（而他们自己的生活却像鱼类一样：因为鱼吞噬一切碰到的东西，强者追逐弱者：事实上，这就是以人肉为食，违背你们和你们祖先为维护一切美善与正义而悉心制定的法律，施行暴行），以致连你们派出的各省总督，都不足以审理那些人对他们提出的控诉；对于这些被控的人，甚至在被击打时，不准不以更多的鞭打奉献自己；在被诋毁时，不准不祝福：因为仅做到公义（公义即是以牙还牙）是不够的，我们还须良善，忍受恶事。\par

% structure: p0150 source=h2 target=section reason=relative-heading-rank; base=h2

\section{第三十五章 基督徒谴责并憎恶一切残忍}

思维健全的人谁会断言，我们既如此品行，却是谋杀者呢？因为我们若非先杀了人，便不能吃人肉。所以，前一指控既是虚假的，若有人就第二项指控问他们是否见过自己所声称的事，没有一个人会如此厚颜无耻地承认见过。然而我们有奴隶，有的多些少些，我们无法不被他们看见；但即便是这些奴隶，也未曾找到一个编造此类事来控告我们。因为他们知道，我们连看到一个人被处死都不能忍受，即使是公正的处死；那么他们中谁能指控我们犯有谋杀或食人罪呢？谁不把角斗士和野兽的搏斗视为最感兴趣的事，尤其是你们所举办的那些？但我们却认为，看到一个人被处死无异于亲手杀死他，因此弃绝了这类表演。那么，我们既连观看都不愿，以免沾染罪责和玷污，又怎能去杀人呢？而且我们说，那些用药堕胎的妇女犯了谋杀罪，必将为堕胎向天主交账，那么我们凭什么去犯谋杀罪呢？因为一个人不会既将胎中的胎儿视为受造物，因而看为天主所眷顾的对象，又在它出世后将它杀掉；既因遗弃婴孩的人犯了杀婴罪而不去遗弃婴孩，却又在养育之后将它毁灭。我们在一切事上始终是表里如一，顺服理性，不凌驾于理性之上。\par

% structure: p0152 source=h2 target=section reason=relative-heading-rank; base=h2

\section{三十六、复活教义对基督徒行为的影响}

谁，若相信复活，会把自己变成那些将要复活的尸体的坟墓呢？因为既相信我们的身体将要复活，却又好像不会复活一样去吃它们，并不属于同一群人的作派；并认为大地会将它所持有的身体交还，但一个人埋藏在自己里面的身体却不会被索还。相反，可以合理地设想，那些认为不必为今生无论善恶度过而交账、没有复活、且认为灵魂与身体一同消亡、好像在其中熄灭的人，将无所不胆大妄为；而那些确信没有事能逃脱天主的鉴察，甚至那为灵魂非理性的冲动及其欲望服务的身体，也将与灵魂一同受罚的人，是不太可能犯最小的罪的。但若有人认为那已经腐朽、分解、归于虚无的身体能被重新构建起来完全是胡说八道，那么对那些不信的人，我们当然不能被合理地指控为邪恶，只能说是愚蠢；因为用来欺骗自己的这些意见，并不伤害他人。然而，并非只有我们相信身体将会复活，许多哲学家也持相同观点，这一点现在不宜展开说明，免得我们被以为引入与当前主题无关的话题，比如谈论\OriginalTerm{可理解者}{intelligible}和\OriginalTerm{可感者}{sensible}以及它们各自的本性，或者争论\OriginalTerm{非形体}{incorporeal}比\OriginalTerm{形体}{corporeal}更古老，\OriginalTerm{可理解者}{intelligible}先于\OriginalTerm{可感者}{sensible}，尽管我们最初熟悉的是后者，因为\OriginalTerm{形体}{corporeal}是由\OriginalTerm{非形体}{incorporeal}与\OriginalTerm{可理解者}{intelligible}结合而成，\OriginalTerm{可感者}{sensible}则是由\OriginalTerm{可理解者}{intelligible}形成的；因为按照\Person{毕达哥拉斯}和\Person{柏拉图}，当身体分解时，它们完全可以从最初构成它们的那些相同元素重新构建起来，这并无阻碍。但让我们暂且搁置关于复活的讨论。\par

% structure: p0154 source=h2 target=section reason=relative-heading-rank; base=h2

\section{三十七、请求公正评判}

如今，您这位在一切事上、天性如此、教养如此、正直、温和、仁厚、配得统治的人，在我清除了各项指控，证明我们虔诚、温和、心灵节制之后，请垂首表示赞同。因为谁比我们这些为您的政权祈祷的人更配得他们所祈求的呢？我们祈求您能按照最公义的方式，子承父业，承受王国，并使您的帝国日益扩展加增，万民都臣服于您的统治之下。这对我们也有益处，使我们能度和平安宁的生活，自己也能欣然遵行一切命令。\par

\SourceRecord{Translated by B.P. Pratten.；Ante-Nicene Fathers；Vol. 2.；Edited by Alexander Roberts, James Donaldson, and A. Cleveland Coxe.；Buffalo, NY: Christian Literature Publishing Co.,；1885.；<http://www.newadvent.org/fathers/0205.htm>.}
